% 第六章 总结与展望
% ch6-conclusion.tex

\subsection{研究总结}

本文围绕奥运项目评估问题，构建了基于AHP-EWM混合权重的决策支持系统。主要研究成果包括：

\subsubsection{评价指标体系}

建立了包含流行度、性别平等、可持续性、包容性、创新性、安全性的六维评价指标体系，系统反映奥运项目评估的多维属性。

\subsubsection{混合权重模型}

创新性地采用两层AHP-EWM混合模型：
\begin{itemize}
\item 第一层：AHP-EWM混合计算六维指标权重
\item 第二层：纯EWM计算子指标权重
\end{itemize}

该模型综合主观专家判断与客观数据信息，提高了决策科学性。

\subsubsection{模型验证}

通过2000-2024年历史数据验证，混合模型在预测准确率上优于单一权重方法。灵敏度分析证明模型结果稳定可靠。

\subsubsection{系统实现}

开发了基于Python的评估系统，提供权重计算、项目评分、可视化展示等功能，为实际决策提供工具支持。

\subsection{主要创新点}

\begin{enumerate}
\item 提出两层AHP-EWM混合权重模型，适应不同数据情况
\item 构建面向奥运项目评估的六维指标体系
\item 实现专家访谈与数据驱动相结合的权重确定方法
\item 开发交互式评估系统，提高决策透明度
\end{enumerate}

\subsection{研究局限性}

\subsubsection{数据局限}

\begin{enumerate}
\item 部分指标数据获取困难（如碳排放、受伤率）
\item 专家访谈样本量有限
\item 历史数据可能不反映未来趋势
\end{enumerate}

\subsubsection{方法局限}

\begin{enumerate}
\item 线性组合模型假设较简单
\item 未考虑指标间的交互作用
\item 预测模型缺乏时间序列分析
\end{enumerate}

\subsection{未来研究方向}

\subsubsection{方法改进}

\begin{enumerate}
\item 引入非线性组合方法
\item 考虑指标相关性修正
\item 结合机器学习提升预测精度
\end{enumerate}

\subsubsection{数据完善}

\begin{enumerate}
\item 扩大专家访谈样本
\item 建立奥运项目数据库
\item 集成实时数据接口
\end{enumerate}

\subsubsection{应用拓展}

\begin{enumerate}
\item 推广至其他体育赛事评估
\item 为各国奥委会提供决策工具
\item 支持残奥项目评估
\end{enumerate}

\subsection{结语}

奥运项目评估是一个复杂的多准则决策问题。本研究通过AHP-EWM混合模型，为这一问题提供了科学、可量化的解决方案。希望本研究能为IOC决策提供参考，为体育管理研究贡献方法论，为后续研究奠定基础。